% !TEX root = main.tex
\section{Related work}
% Each prior-work sentence carries its own gap. No claim that this paper studies language models.

\paragraph{Description length.}
\citet{kolmogorov1968three} measure information by the shortest program that produces an object on a fixed machine.
\citet{grunwald2008algorithmic} separate that description length from the time the program takes to run.
\citet{rissanen1978modeling} selects models by total description length, charging for the model as well as the data.
Final cross-entropy under a trained network omits that charge.
\citet{voita2020information} make the charge operational by coding labels online, before each update of a fixed learning rule.
Our online code in \texttt{experiments/archive/07\_prequential\_mdl\_and\_length\_transfer/} follows that prequential pattern on the addition stream.
It is not an estimate of Kolmogorov complexity.
The conditional complexity of each test pair stays fixed because the input, the answer, the task statement, and the reference machine stay fixed.

\paragraph{What a ranking is attached to.}
\citet{bengio2009curriculum} train on easier examples before harder ones, using a difficulty signal chosen in advance.
That design is useful when the order is stable.
Our measurements ask whether the order stays put after the mixture changes.
\citet{lake2018generalization} find that sequence models can fit a compositional training set and still miss systematic recombinations.
\citet{keysers2020measuring} separate a match of primitive distributions from a match of their compositions.
Those papers study sequence tasks.
Our carry design is a smaller arithmetic version of the same question, with an exact checker and a frozen test set.
\citet{zhang2017understanding} show that deep networks can fit random labels of a finite training set.
Fitting those labels does not predict independent random labels on new items.
Our shuffled-label control is the addition analogue: exact-sum accuracy stays near zero.

\paragraph{Resource laws and metric artifacts.}
\citet{kaplan2020scaling} and \citet{hoffmann2022chinchilla} predict language-model loss from scale and from data or training compute.
Those laws describe average loss.
They do not identify which evaluated items were in the training mixture.
We borrow only the idea of a simple law in parameter count and updates, then test extra covariates on addition networks.
This is not a scaling-law result for language models.
\citet{power2022grokking} show that, on small algorithmic tasks, generalization can arrive long after the training loss has fallen.
Their intervention is training time under one data distribution.
Ours holds the scored budget fixed inside a comparison and changes the mixture.
\citet{schaeffer2023mirage} show that an aggregate metric can create an apparent jump in capability.
That paper is about metric discontinuities on language-model benchmarks.
This paper is about training history on one addition task.
We cite it only as a neighboring warning against reading a score as a property of the task.

\paragraph{Group-relative updates.}
\citet{shao2024deepseekmath} introduce group relative policy optimization for language-model mathematics.
We use the name only for the leave-one-out advantage in our protocol, with group size 16 and a 0/1 exact-sum reward.
We do not train a language model, and we do not reproduce that paper's experiments.
The phrase in our last study means the reward is the exact-sum check.
A cleared public source for that phrase is still missing \todocite{public source for the name reinforcement learning with verifiable rewards}.
A cleared citation for a companion mathematics grid used in the second study is also unresolved.
That grid is described from the project report, not cited as a peer-reviewed law.
